\worksheettitle
  {Названия и диагонали}
  {\modulemark{M06} Версия учителя}

\begin{teacherbox}[Паспорт модуля]
  \textbf{Результат:} ученик проверяет имя многоугольника последовательным
  обходом, различает сторону и диагональ, проводит диагонали из выбранной
  вершины и читает границы полученных частей.

  \textbf{Предварительно:} различает вершины и стороны многоугольника.
\end{teacherbox}

\begin{teacherbox}[Методический акцент]
  При ошибке в названии просите провести пальцем путь, заданный буквами.
  Формулу числа диагоналей не вводите до достаточной практики с рисунками;
  она относится к M06\(\star\).
\end{teacherbox}

\section{Вход и разобранный пример}

Во входе: соседняя с $A$ вершина $B$; имя $ACBD$ неверно; проверка ---
последовательный обход сторон.

\correctnamingfigure

\begin{practicebox}[1. Проверка имён]
  Верны $CDEFAB$, $FEDCBA$, $DEFCBA$. Запись $ABDEFC$ неверна:
  первый нарушенный переход $B-D$.
\end{practicebox}

\section{Стороны и диагонали}

\diagonalfigure

В пропусках: соседние; несоседние; $AC$, $AD$, $AE$.

\begin{practicebox}[2. Диагонали из $A$]
  \hexagondiagonalsanswer
  Диагонали $AC$, $AD$, $AE$; всего три.
\end{practicebox}

\begin{practicebox}[3. Классификация]
  $AB$ --- сторона; $AC$ --- диагональ; $FA$ --- сторона;
  $CE$ --- диагональ.
\end{practicebox}

\pentagonsplitfigure

\begin{practicebox}[4. Полученные части]
  Треугольник $ABC$ и четырёхугольник $ACDE$. Допустимы циклические
  переименования и обратный порядок обхода.
\end{practicebox}

\section{Ошибка и контроль}

\begin{warningbox}[Ожидаемое исправление]
  Причина: не проверено соседство вершин. $B$ и $C$ соседние, поэтому $BC$
  является стороной. Диагональ соединяет несоседние вершины, например $BD$.
\end{warningbox}

\begin{checkpointbox}[Ответы]
  \begin{enumerate}[label=\textbf{\arabic*.}]
    \item Верны $ABCDEF$, $CDEFAB$, $AFEDCB$; $AEDCBF$ нарушает обход.
    \item Из $C$: $CA$, $CE$, $CF$.
    \item Допустимы обозначения $AC$, $CE$, $CF$.
    \item $C$ и $D$ соседние вершины, поэтому $CD$ --- сторона.
  \end{enumerate}
\end{checkpointbox}

\begin{teacherbox}[Решение о маршруте]
  M06-R назначается, если имя не связывается с обходом или соседство не
  проверяется. M06-H достаточно при верном способе и единичном пропуске.
  M06\(\star\) предлагается после устойчивого базового контроля.
\end{teacherbox}
